\documentclass[xcolor=table,aspectratio=169]{beamer}

%\usepackage[english]{babel} % English language/hyphenation

\usepackage[defaultsans,scale=0.95]{opensans}
\usepackage{FiraMono}
\usepackage{fourier} % Use the Adobe Utopia font for the document - comment this line to return to the LaTeX default
\usefonttheme{professionalfonts} % https://tex.stackexchange.com/a/330700/235983
\DeclareMathAlphabet{\mathcal}{OMS}{cmsy}{m}{n}
\newcommand\bmmax{2} % used to fix issues with bbm and e.g. esint
\usepackage{bm} % for vec/mat
\usepackage[T1]{fontenc}

\usepackage{graphicx}
\usepackage{caption}
\captionsetup[figure]{labelformat=empty}
\usepackage{subcaption}
\usepackage{multicol}

\usepackage{tikz}
\usetikzlibrary{decorations.pathreplacing}
\usepackage{wrapfig}

% COLORS
\usepackage{xcolor}
\colorlet{myred}{red!80!black}
\colorlet{myblue}{blue!80!black}
\colorlet{mygreen}{green!60!black}
\colorlet{myorange}{orange!70!red!60!black}
\colorlet{mydarkred}{red!30!black}
\colorlet{mydarkblue}{blue!40!black}
\colorlet{mydarkgreen}{green!30!black}

\colorlet{mylightblue}{blue!60!white}
\colorlet{mydarkyellow}{yellow!70!black}
\colorlet{myyellowred}{yellow!70!red}

\usepackage[backend=biber,style=numeric,sorting=none]{biblatex}
\addbibresource{template_for_qibin_bib.bib}

%\usepackage{tabu}  % \(\)-able tabulars
\usepackage{float} % H placement
\usepackage{hyperref}

\usepackage{amsmath, amsthm, amssymb, amsfonts}
% \usepackage{mbboard} % requires manual installation in texmf
\usepackage{esint}
\usepackage{mathtools} % boxed aligns, delimiters
\usepackage{listings}
\lstset{basicstyle=\ttfamily}

\usepackage{physics} % http://mirrors.ibiblio.org/CTAN/macros/latex/contrib/physics/physics.pdf
\usepackage[makeroom]{cancel}
% \usepackage{slashed}
\usepackage{siunitx}
\sisetup{inter-unit-product=\ensuremath{{}\cdot{}}}
\DeclareSIUnit\eVperc{\eV\per\clight}
\AtBeginDocument{\RenewCommandCopy\qty\SI}
\usepackage[version=4]{mhchem}
%\usepackage{hepnames}
\usepackage[italic]{heppennames}
\usepackage{xpatch} % https://tex.stackexchange.com/a/541095
\makeatletter
\xpatchcmd\@HepConStyle
{\edef\@upcode{\updefault}}
{\ifdefined\shapedefault\edef\@upcode{\shapedefault}\else\edef\@upcode{\updefault}\fi}
{}{}
\makeatother
\usepackage{tensor}
\usepackage{feynmp-auto}
\usepackage{listofitems}
\usetikzlibrary{arrows.meta} % for arrow size
\usepackage[outline]{contour} % glow around text

\addtobeamertemplate{frametitle}{}{\vspace{-15.5pt}}

%%%%%%%%%%%%%%%%%%%%%%%%%%%%%%%%%%%%%%

\newcommand{\fixme}{\text{{\Huge \(\mathbb{FIXME}\)}}}
\newcommand{\todo}{\text{{\huge \(\mathbb{TODO}\)}}}
\newcommand{\lboxed}[1]{\begin{array}{|l}
		\hline\raisebox{3pt}{\(\displaystyle\vphantom{#1}\)}
		\displaystyle#1\!\!\!\!\!
		\raisebox{-3pt}{\(\displaystyle\vphantom{#1}\)}\\\hline
	\end{array}
}
\newcommand{\rboxed}[1]{\begin{array}{r|}
		\hline\raisebox{3pt}{\(\displaystyle\vphantom{#1}\)}
		\!\!\!\!\!\displaystyle#1
		\raisebox{-3pt}{\(\displaystyle\vphantom{#1}\)}\\\hline
	\end{array}
}
\newcommand{\numberthis}{\addtocounter{equation}{1}\tag{\theequation}}% https://tex.stackexchange.com/a/42728

\newcommand{\Lagr}{\mathcal{L}}

\setlength\arrayrulewidth{1.0pt}

\usepackage{booktabs}
\usepackage{subcaption}
%\newcommand\blfootnote[1]{%
%	\begingroup
%	\renewcommand\thefootnote{}\footnote{#1}%
%	\addtocounter{footnote}{-1}%
%	\endgroup
%}
% Macro to skip item in enum
% https://tex.stackexchange.com/a/201162/235983
\newcommand{\skipitems}[1]{%
	\addtocounter{\@enumctr}{#1}%
}
% Macros to start and end Backup section with unnumbered slides
\newcommand{\backupbegin}{
    \newcounter{finalframe}
    \setcounter{finalframe}{\value{framenumber}}
    \setbeamertemplate{footline}{} % Remove the footline for Backup section
    \frame{\usebeamertemplate*{backup page}}
}

\newcommand{\backupend}{
    \setcounter{framenumber}{\value{finalframe}}
    \setbeamertemplate{footline}[frame number] % Restore original slide numbering
}
% Macro to add TOC entry for starred section
% https://tex.stackexchange.com/a/611664/235983
\makeatletter
\newcommand{\nonumsec}{%
	\addtocontents{toc}{\protect\beamer@sectionintoc{\the\c@section}{\secname}{\the\c@page}{\the\c@part}{0}}%
}
\makeatother

\def\thisframelogos{}

\newcommand{\framelogo}[1]{\def\thisframelogos{#1}}

\addtobeamertemplate{frametitle}{}{%
\begin{tikzpicture}[remember picture,overlay]
\node[anchor=north east,yshift=3.5pt] at (current page.north east) {%
    \foreach \img in \thisframelogos {%
        \hspace{.5ex}%
        \includegraphics[height=0.8cm, keepaspectratio]{\img}%
    }%
};
\end{tikzpicture}}

%%%%%%%%%%%%%%%%%%%%%%%%%%%%%%%%%%%%%%

\usetheme{Philadelphia}
%\usefonttheme[onlymath]{serif}
\hypersetup{colorlinks,allcolors=pennred}

\title[ACAT 2025]{Deep Learning for Primary Vertex Identification \\ in the ATLAS Experiment}

\author[Qi Bin Lei]{\href{https://www.linkedin.com/in/qibinlei/}{Qi Bin~Lei}, Rocky B. Garg, Lauren Tompkins}

\institute[]{Stanford University}

\date{September 11th, 2025}

% Replace these with the path to the logos
\titlegraphic{
    \hspace{5pt}
	\href{https://stanford.edu/group/stanford_atlas/}{\includegraphics[height=1cm]{logo-Stanford/Block_S_2_color.png}}
	\href{https://atlas.cern/}{\includegraphics[trim={0 11.5cm 0 11.5cm},clip,height=1cm]{logo-ATLAS/ATLAS-Logo-Ref-Pantone.pdf}}
    \href{https://iris-hep.org/}{\includegraphics[height=1cm]{logo-iris-hep/Iris-hep-4-no-long-name.png}}
}

\begin{document}

% \begin{frame}{Title}
% 	\begin{itemize}
% 		\item An item
% 	\end{itemize}
% 	\begin{columns}
% 		\column{0.5\linewidth}
% 		Hello
% 		\column{0.5\linewidth}
% 		World
% 	\end{columns}
% \end{frame}
% \section{Section 1}
% \frame{\sectionpage}

\framelogo{"logo-ATLAS/ATLAS-Logo-White-transparent.png"}

\begin{frame}[plain]
    \titlepage
\end{frame}

\begin{frame}{Current Primary Vertex Algorithm: AVMF}
    % Talk about innate challenges that comes with using AVMF
    \begin{columns}
        % Text
        \column{0.4\linewidth}
        \begin{itemize}
            \item Global approach to vertex finding and fitting
            \item Finds and fits tracks based on compatibility to different primary vertices
            \item Limitations
            \begin{enumerate}
                \item Parallelization
                \item Time scale increases non linearly with number of PVs
            \end{enumerate}
        \end{itemize}

        % Text
        \column{0.6\linewidth}
        \begin{figure}
            \centering
            \includegraphics[width=0.9\linewidth]{oct25th/AVMF.png}
            \label{fig:AVMF}
            \caption{\href{https://cds.cern.ch/record/2670380/files/ATL-PHYS-PUB-2019-015.pdf}{[ATL-PHYS-PUB-2019-015]}}
        \end{figure}
    \end{columns}
\end{frame}

\begin{frame}{PV-Finder Before}
    \vspace{5pt}
    \centering
    This network performs nicely but is limited by execution time
    \vspace{25pt}
    \begin{figure}
        \centering
        \includegraphics[width=1\linewidth]{APS2025/analyticalkde_pv.png}
        \label{fig:oldPVFinderArchitecture}
    \end{figure}
\end{frame}

\begin{frame}{PV-Finder After}
    \centering
    TVPD presentation on PV-Finder results has more details! (\href{https://indico.cern.ch/event/1534661/}{TVPD, April 4th})
    \begin{figure}
        \centering
        \includegraphics[width=1\linewidth]{april4th/updated_tracks_to_pv.png}
        \label{fig:oldPVFinderArchitecture}
    \end{figure}
\end{frame}

\begin{frame}{Example End-to-end tracks-to-hist Output}
    \begin{figure}
        \centering
        \includegraphics[width=0.8\linewidth]{june25th/example_end_to_end_with_class.png}
        \label{fig:example_Hists_output_1}
    \end{figure}
\end{frame}

\begin{frame}{\(\Delta z\) Resolution and \(\sigma_{\text{vtx-vtx}}\) Calculation}
    \begin{columns}
        \begin{column}{0.6\linewidth}
            \begin{figure}
                \centering
                \includegraphics[width=0.9\linewidth]{april4th/deltaz(2).png}
            \end{figure}
        \end{column}

        \begin{column}{0.4\linewidth}
            \begin{figure}
                \centering
                \includegraphics[width=1\linewidth]{april4th/sigma_vtx_vtx_dropout.png}
            \end{figure}
        \end{column}
    \end{columns}
\end{frame}

\begin{frame}{UNet vs UNet++ Efficiency and FPR with MLP}
    \centering
    UNet and UNet++ perform basically identically!
    \begin{figure}
        \centering
        \includegraphics[width=1\linewidth]{april4th/droput_unet_unetplusplus_eff.png}
    \end{figure}
\end{frame}

\begin{frame}{Intermission}
    \begin{itemize}
        \item Great that PV-Finder can reconstruct more clean/merged PV!
        \item Tracks to Vertex Association?
        \begin{itemize}
            \item Necessary for hard-scatter vertex calculations
            \item Allows for proper comparison with AMVF
        \end{itemize}
        \item \textbf{GNN for TTVA}
        \begin{itemize}
            \item Graph Construction
            \item AMVF-Stype Vertex Classification
            \item Preliminary Performance
            \item Case Study
        \end{itemize}
    \end{itemize}
\end{frame}

\begin{frame}{PV-Finder + GNN Fitter}
    \begin{figure}
        \centering
        \includegraphics[width=0.55\linewidth]{april4th/pipeline_pvfinder_gnn.png}
    \end{figure}
\end{frame}

\begin{frame}{GNN Concept}
    \begin{figure}
        \centering
        \includegraphics[width=0.55\linewidth]{april4th/gnn_concept.png}
    \end{figure}
\end{frame}

\begin{frame}{Forming Graphs}
    \begin{figure}
        \centering
        \includegraphics[width=0.6\linewidth]{april4th/forming_edges_gnn.png}
    \end{figure}
\end{frame}

\begin{frame}{Scoring Edges in Graphs}
    \begin{figure}
        \centering
        \includegraphics[width=0.6\linewidth]{april4th/scoring_edges_gnn.png}
    \end{figure}
\end{frame}

\begin{frame}{Choosing Edges in Graphs}
    \begin{figure}
        \centering
        \includegraphics[width=0.6\linewidth]{april4th/choosing_edges_gnn.png}
    \end{figure}
\end{frame}

\begin{frame}{GNN Pipeline}
    \begin{figure}
        \centering
        \includegraphics[width=0.9\linewidth]{april15th/endtoend_gnn_pipeline.png}
    \end{figure}
\end{frame}

\begin{frame}[plain]
    \centering
    AMVF-Style Vertex Classification
\end{frame}

\begin{frame}{Vertex Classification (1)}
    \begin{figure}
        \centering
        \includegraphics[width=0.7\linewidth]{june20th/reco_vertex_classification_1.png}
    \end{figure}
\end{frame}

\begin{frame}{Vertex Classification (2)}
    \begin{figure}
        \centering
        \includegraphics[width=0.7\linewidth]{june20th/reco_vertex_classification_2.png}
    \end{figure}
\end{frame}

\begin{frame}{Vertex Classification (3)}
    \begin{figure}
        \centering
        \includegraphics[width=0.7\linewidth]{june20th/reco_vertex_classification_3.png}
    \end{figure}
\end{frame}

\begin{frame}{Vertex Classification (4)}
    \begin{figure}
        \centering
        \includegraphics[width=0.7\linewidth]{june20th/reco_vertex_classification_4.png}
    \end{figure}
\end{frame}

\begin{frame}{Vertex Classification (5)}
    \begin{figure}
        \centering
        \includegraphics[width=0.7\linewidth]{june20th/reco_vertex_classification_5.png}
    \end{figure}
\end{frame}

\begin{frame}{Vertex Classification (6)}
    \begin{figure}
        \centering
        \includegraphics[width=0.7\linewidth]{june20th/reco_vertex_classification_6.png}
    \end{figure}
\end{frame}

\begin{frame}{Vertex Classification (7)}
    \begin{figure}
        \centering
        \includegraphics[width=0.7\linewidth]{june20th/reco_vertex_classification_7.png}
    \end{figure}
\end{frame}

\begin{frame}{Vertex Classification (8)}
    \begin{figure}
        \centering
        \includegraphics[width=0.7\linewidth]{june25th/reco_vertex_classification_8.png}
    \end{figure}
\end{frame}

\begin{frame}[plain]
    \centering
    Evaluating Reco Vertices with AMVF-Style Vertex Classification
\end{frame}

\begin{frame}{Steps Outlined}
    \begin{itemize}
        \item Train PV-Finder
        \begin{itemize}
            \item Using current best PV-Finder run (UNet with Dropout Model)
            \item (Old Metrics) Efficiency: 94.56\%
            \item (Old Metrics) FPR/Average Number of Fake Reco PV: 0.72
        \end{itemize}
        \item Calculate predicted primary vertex location
        \item Construct graphs from reconstructed PV information
        \item Evaluate using GNN model trained on truth information
        \begin{itemize}
            \item Using current best GNN run (GraphConv Model)
            \item Achieves results shown in prior slides
        \end{itemize}
    \end{itemize}
\end{frame}

% \begin{frame}{Re: Changes and Methods for Evaluation}
%     \begin{itemize}
%         \item Every track will now have a potential edge with every PV
%         \begin{itemize}
%             \item \textbf{No more} \( \abs{z_{0, \text{track}} - z_{\text{pv}}} <= 5\text{mm} \text{ or } 10\text{mm}\)
%         \end{itemize}
%         \item Outputs of the GNN can be interpreted as weights and evaluated in many ways
%         \item Method 1 (Threshold): Threshold of having a score be greater than 0.7, 0.85, or 0.9
%         \item Method 2 (MaxScore): Highest scoring edge for each track is 1, rest is 0
%     \end{itemize}
% \end{frame}

\begin{frame}{Reco Vertices with Classification at \(\left< \mu \right> =60\) (Thr=0.95)}
    \begin{columns}
        \begin{column}{0.5\textwidth}
            \begin{itemize}
                \item Total Reconstructed PVs: 62457
                \item Total Truth PVs: 80498
                \item Total Clean: 42880, Rate: 68.66\%
                \item Total Merged: 13437, Rate: 21.51\%
                \item Total Split: 3838, Rate: 9.83\%
                \item Total Fake: 3, Rate: 0\%
            \end{itemize}
        \end{column}

        \begin{column}{0.5\textwidth}
            \begin{figure}
                \centering
                \includegraphics[width=1\linewidth]{june25th/reco_pv_avg_bar_chart_95thr_recovertex.png}
                \caption{Number of reco vertices at \(\left< \mu\right> = 60\)}
            \end{figure}
        \end{column}
    \end{columns}
\end{frame}

% \begin{frame}{Reco Vertices with Classification at \(\left< \mu \right> =60\) MaxScore}
%     \begin{columns}
%         \begin{column}{0.55\textwidth}
%             \begin{itemize}
%                 \item Total Reconstructed PVs: 62457
%                 \item Total Clean: 49500, Rate: 0.7925
%                 \item Total Merged: 8243, Rate: 0.1320
%                 \item Total Split: 3838, Rate: 0.0615
%                 \item Total Fake: 876, Rate: 0.0140
%             \end{itemize}
%         \end{column}

%         \begin{column}{0.45\textwidth}
%             \begin{figure}
%                 \centering
%                 \includegraphics[width=1\linewidth]{example-image-a}
%             \end{figure}
%         \end{column}
%     \end{columns}
% \end{frame}

\begin{frame}{Comparison of Reco Vertex Classification: Thresholds (1)}
    \begin{columns}
        \begin{column}{0.5\textwidth}
            \begin{figure}
                \centering
                \includegraphics[width=1\linewidth]{june25th/stacked_pv_classification_5_80_85thr_recovertex.png}
                \caption{Thr = 0.85}
            \end{figure}
        \end{column}

        \begin{column}{0.5\textwidth}
            \begin{figure}
                \centering
                \includegraphics[width=1\linewidth]{june25th/stacked_pv_classification_5_80_95thr_recovertex.png}
                \caption{Thr = 0.95}
            \end{figure}
        \end{column}
    \end{columns}
\end{frame}

\begin{frame}{Comparison of Reco Vertex Classification: Thresholds (2)}
    \begin{columns}
        \begin{column}{0.5\textwidth}
            \begin{figure}
                \centering
                \includegraphics[width=1\linewidth]{june5th/Screenshot 2025-06-10 at 11.33.19 PM.png}
                \caption{AMVF Pub Note: \href{https://cds.cern.ch/record/2670380/files/ATL-PHYS-PUB-2019-015.pdf}{Link}}
            \end{figure}
        \end{column}

        \begin{column}{0.5\textwidth}
            \begin{figure}
                \centering
                \includegraphics[width=1\linewidth]{june25th/stacked_pv_classification_5_80_95thr_recovertex.png}
                \caption{Thr = 0.95}
            \end{figure}
        \end{column}
    \end{columns}
\end{frame}

% \begin{frame}{Comparison of Reco Vertex Classification: MaxScore}
%     \begin{columns}
%         \begin{column}{0.5\textwidth}
%             \begin{figure}
%                 \centering
%                 \includegraphics[width=1\linewidth]{example-image-a}
%                 \caption{MaxScore}
%             \end{figure}
%         \end{column}

%         \begin{column}{0.5\textwidth}
%             \begin{figure}
%                 \centering
%                 \includegraphics[width=1\linewidth]{june5th/stacked_pv_classification_5_80_thr95_recovertex.png}
%                 \caption{Thr = 0.95}
%             \end{figure}
%         \end{column}
%     \end{columns}
% \end{frame}

\begin{frame}{A Look into GNN Associations and Thresholds}
    \begin{itemize}
        \item Classification rates skew less clean with AMVF-style definition but are similar
        \item How do the track-vertex associations look?
        \item What is causing the increase in merged vertices?
        \item Are there tracks that are not being used due to threshold?
        \item Lets look at test cases!
    \end{itemize}
    \begin{figure}
        \centering
        \includegraphics[width=0.4\linewidth]{this_is_fine.jpg}
    \end{figure}
\end{frame}

\begin{frame}{Case Study: Looking at Merged/Split Reco PVs}
    \begin{columns}
        \begin{column}{0.5\textwidth}
            \begin{itemize}
                \item Total Reco PVs: 42
                \item Total Truth PVs: 59
                \item Thr=0.85
                \begin{itemize}
                    \item Num Clean: 22
                    \item Num Merged: 12
                    \item Num Split: 8
                    \item Num Fake: 0
                \end{itemize}
                \item Thr=0.95
                \begin{itemize}
                    \item Num Clean: 30
                    \item Num Merged: 7
                    \item Num Split: 5
                    \item Num Fake: 0
                \end{itemize}
            \end{itemize}
        \end{column}

        \begin{column}{0.5\textwidth}
            \begin{figure}
                \centering
                \includegraphics[width=1\linewidth]{june25th/casestudy_example_reco_truth_hist_dist_31_38_51.png}
            \end{figure}
        \end{column}
    \end{columns}
\end{frame}

\begin{frame}{Case Study: Three Nearby Vertices Thr=0.85 (1)}
    \begin{columns}
        \begin{column}{0.4\textwidth}
            \begin{itemize}
                \item Reco PV 24: 33 Assoc tracks
                \begin{itemize}
                    \item Old Class: Merged
                    \item \(W_{\text{reco},24} = 33\)
                    \item \(\sum p_T^2 = 4e7\) (GeV\(^2\))
                \end{itemize}
                \item Truth PV 31: 13 Assoc tracks
                \begin{itemize}
                    \item \(W_{\text{truth},31} = 12\)
                    \item \(W_{\text{truth},31}/W_{\text{reco},24} = 0.4\)
                \end{itemize}
                \item Truth PV 38: 10 Assoc tracks
                \begin{itemize}
                    \item \(W_{\text{truth},38} = 10\)
                    \item \(W_{\text{truth},38}/W_{\text{reco},24} = 0.3\)
                \end{itemize}
                \item Truth PV 51: 6 Assoc tracks
                \begin{itemize}
                    \item \(W_{\text{truth},51} = 5\)
                    \item \(W_{\text{truth},51}/W_{\text{reco},24} = 0.15\)
                \end{itemize}
                \item AMVF-style Class: \textbf{Split}
            \end{itemize}
        \end{column}

        \begin{column}{0.6\textwidth}
            \begin{figure}
                \centering
                \includegraphics[width=1\linewidth]{june25th/casestudy_example_reco_truth_hist_dist_24.png}
            \end{figure}
        \end{column}
    \end{columns}
\end{frame}

\begin{frame}{Case Study: Three Nearby Vertices Thr=0.85 (2)}
    \begin{columns}
        \begin{column}{0.4\textwidth}
            \begin{itemize}
                \item Reco PV 25: 42 Assoc tracks
                \begin{itemize}
                    \item Old Class: Clean
                    \item \(\sum p_T^2 = 6e7\) (GeV\(^2\))
                    \item \(W_{\text{reco},25} = 42\)
                \end{itemize}
                \item Truth PV 31: 13 Assoc tracks
                \begin{itemize}
                    \item \(W_{\text{truth},31} = 12\)
                    \item \(W_{\text{truth},31}/W_{\text{reco},25} = 0.4\)
                \end{itemize}
                \item Truth PV 38: 10 Assoc tracks
                \begin{itemize}
                    \item \(W_{\text{truth},38} = 10\)
                    \item \(W_{\text{truth},38}/W_{\text{reco},25} = 0.3\)
                \end{itemize}
                \item Truth PV 51: 6 Assoc tracks
                \begin{itemize}
                    \item \(W_{\text{truth},51} = 5\)
                    \item \(W_{\text{truth},51}/W_{\text{reco},25} = 0.15\)
                \end{itemize}
                \item AMVF-style Class: \textbf{Merged}
            \end{itemize}
        \end{column}

        \begin{column}{0.6\textwidth}
            \begin{figure}
                \centering
                \includegraphics[width=1\linewidth]{june25th/casestudy_example_reco_truth_hist_dist_25.png}
            \end{figure}
        \end{column}
    \end{columns}
\end{frame}

\begin{frame}{Case Study: Three Nearby Vertices Thr=0.95 (1)}
    \begin{columns}
        \begin{column}{0.4\textwidth}
            \begin{itemize}
                \item Reco PV 24: 24 Assoc tracks
                \begin{itemize}
                    \item \(W_{\text{reco},24} = 24\)
                    \item \(\sum p_T^2 = 4e7\) (GeV\(^2\))
                \end{itemize}
                \item Truth PV 31: 13 Assoc tracks
                \begin{itemize}
                    \item \(W_{\text{truth},31} = 12\)
                    \item \(W_{\text{truth},31}/W_{\text{reco},24} = 0.5\)
                \end{itemize}
                \item Truth PV 38: 10 Assoc tracks
                \begin{itemize}
                    \item \(W_{\text{truth},38} = 6\)
                    \item \(W_{\text{truth},38}/W_{\text{reco},24} = 0.25\)
                \end{itemize}
                \item Truth PV 51: 6 Assoc tracks
                \begin{itemize}
                    \item \(W_{\text{truth},51} = 5\)
                    \item \(W_{\text{truth},51}/W_{\text{reco},24} = 0.21\)
                \end{itemize}
                \item AMVF-style Class: \textbf{Merged}
            \end{itemize}
        \end{column}

        \begin{column}{0.6\textwidth}
            \begin{figure}
                \centering
                \includegraphics[width=1\linewidth]{june25th/casestudy_example_reco_truth_hist_dist_24.png}
            \end{figure}
        \end{column}
    \end{columns}
\end{frame}

\begin{frame}{Case Study: Three Nearby Vertices Thr=0.95 (2)}
    \begin{columns}
        \begin{column}{0.4\textwidth}
            \begin{itemize}
                \item Reco PV 25: 22 Assoc tracks
                \begin{itemize}
                    \item \(W_{\text{reco},25} = 22\)
                    \item \(\sum p_T^2 = 3e7\) (GeV\(^2\))
                \end{itemize}
                \item Truth PV 31: 13 Assoc tracks
                \begin{itemize}
                    \item \(W_{\text{truth},31} = 9\)
                    \item \(W_{\text{truth},31}/W_{\text{reco},25} = 0.4\)
                \end{itemize}
                \item Truth PV 38: 10 Assoc tracks
                \begin{itemize}
                    \item \(W_{\text{truth},38} = 9\)
                    \item \(W_{\text{truth},38}/W_{\text{reco},25} = 0.4\)
                \end{itemize}
                \item Truth PV 51: 6 Assoc tracks
                \begin{itemize}
                    \item \(W_{\text{truth},51} = 0\)
                \end{itemize}
                \item AMVF-style Class: \textbf{Merged}
            \end{itemize}
        \end{column}

        \begin{column}{0.6\textwidth}
            \begin{figure}
                \centering
                \includegraphics[width=1\linewidth]{june25th/casestudy_example_reco_truth_hist_dist_25.png}
            \end{figure}
        \end{column}
    \end{columns}
\end{frame}

\begin{frame}{Number of Associations Per Tracks Varying Threshold}
    \begin{columns}
        \begin{column}{0.5\textwidth}
            \begin{figure}
                \centering
                \includegraphics[width=1\linewidth]{june25th/redo_casestudy_num_track_assocs_thr85.png}
            \end{figure}
        \end{column}

        \begin{column}{0.5\textwidth}
            \begin{figure}
                \centering
                \includegraphics[width=1\linewidth]{june25th/redo_casestudy_num_track_assocs_thr95.png}
            \end{figure}
        \end{column}
    \end{columns}
\end{frame}

\begin{frame}{Re: A Look into GNN Associations and Thresholds}
    \begin{itemize}
        \item What is causing the increase in merged/split vertices?
        \begin{itemize}
            \item Seems to be that distinguishing tracks for nearby vertices are tough
        \end{itemize}
        \item Are there tracks that are not being used due to threshold?
        \begin{itemize}
            \item Tracks are being used even in high thresholds
        \end{itemize}
        \item Are there other ways on evaluating performance?
        \begin{itemize}
            \item \textbf{Hard Scatter Vertex Classification}
        \end{itemize}
    \end{itemize}
    \begin{figure}
        \centering
        \includegraphics[width=0.4\linewidth]{this_is_fine.jpg}
    \end{figure}
\end{frame}

\begin{frame}{Concluding Remarks}
    \begin{itemize}
        \item Presented a ML tool to find and fit primary vertices
        \item How to continue improving performance...? What is the next direction...?
        \begin{itemize}
            \item \textbf{Modifying Graphs}
            \begin{itemize}
                \item Edge Attributes
                \item Binary Cross Entropy to Cross Entropy
            \end{itemize}
            \item Start applying cuts to the output of GNN
            \item Limit 1 edge per track
            \item \textbf{Using Reco PVs to train GNN}
        \end{itemize}
        \item Developing tools to continue analyzing PV-Finder + GNN outputs
        \begin{itemize}
            \item \textbf{HS Vertex Quality Metrics}
        \end{itemize}
        \item Abstract accepted for ACAT 2025!
    \end{itemize}
\end{frame}

%%%%% Backup %%%%%%
\section*{Backup}
\nonumsec
\backupbegin
% \begin{frame}{Network Structure Readout}
% 	\begin{wrapfig}{r}{0.5\textwidth}
% 	    \centering
% 	    \includegraphics[width=.6\textwidth]{graphs/DenseLayers.png}
%         \caption*{Dense Layers}
% 	\end{wrapfig}
% \end{frame}

\begin{frame}{Relevant Links and Directories}
    \begin{itemize}
        \item LHCb PV-Finder
        \begin{itemize}
            \item \href{https://gitlab.cern.ch/sakar/pv-finder_v2/-/tree/master}{Link to Gitlab}
        \end{itemize}
        \item ATLAS PV-Finder
        \begin{itemize}
            \item \href{https://gitlab.cern.ch/qlei/atlas_pvfinder}{Link to Gitlab}
        \end{itemize}
        \item Prior Internal/Conference Presentations:
        \begin{itemize}
            \item TVPD Oct 2024: \href{https://drive.google.com/file/d/17qCnL57fT7LCKyQ1Rz-soJySaNCVeMRV/view?usp=sharing}{Slides}
            \item US LUA 2024: \href{https://drive.google.com/file/d/1HbAjlKBw6nej-oqOksUL_HOfKY8xsd2r/view?usp=sharing}{Slides}
            \item APS Joint Meeting March 2025: \href{https://drive.google.com/file/d/1j6pnqzuXcP0OWXoUXkwElH8lAbv9GtRq/view?usp=drive_link}{Slides}
            \item TVPD April 2025: \href{https://indico.cern.ch/event/1534661/}{Slides}
            \item TVPD June 2025: \href{https://indico.cern.ch/event/1557481/}{Slides}
        \end{itemize}
    \end{itemize}
\end{frame}

\begin{frame}[plain]
    \centering
    PV-Finder More
\end{frame}

\begin{frame}{A Brief Review of UNet}
    \begin{columns}
        \begin{column}{0.4\linewidth}
            \begin{itemize}
                \item Effective for feature extraction
                \item Downsampling reduces dimensionality, helping emphasize important data features
                \item Skip connections propogate prior information
                \item \(\sim 279,000\) Parameters
            \end{itemize}
        \end{column}

        \begin{column}{0.6\linewidth}
            \begin{figure}
                \centering
                \includegraphics[width=1\linewidth]{april4th/UNet.png}
                \label{fig:pipeline_kdetohist}
            \end{figure}
        \end{column}
    \end{columns}
\end{frame}

\begin{frame}{A Brief Review of UNet++}
    \begin{columns}
        \begin{column}{0.4\linewidth}
            \begin{itemize}
                \item Just more skip connections and upsamplings
                \item \(\sim 796,000\) Parameters
            \end{itemize}
        \end{column}

        \begin{column}{0.6\linewidth}
            \begin{figure}
                \centering
                \includegraphics[width=1\linewidth]{april4th/unetplusplus_figure.png}
            \end{figure}
        \end{column}
    \end{columns}
\end{frame}

\begin{frame}{Comparison with AVMF: Vertex Classification Scheme}
    \begin{figure}
        \centering
        \includegraphics[width=0.75\linewidth]{APS2025/vertex_classification_aps2025.png}
    \end{figure}
\end{frame}

\begin{frame}{Number of Primary Vertives vs \(\mu\)}
    \begin{figure}
        \centering
        \includegraphics[width=0.7\linewidth]{APS2025/npv_compare_Total.png}
    \end{figure}
\end{frame}

\begin{frame}{Vertex Resolution vs Tracks}
    \begin{figure}
        \centering
        \includegraphics[width=0.7\linewidth]{APS2025/vertex_resolution_tracks.pdf}
    \end{figure}
\end{frame}

\begin{frame}{AMVF CPU Time Per Event}
    \begin{figure}
        \centering
        \includegraphics[width=0.7\linewidth]{APS2025/acts_pvfinder_timing.png}
    \end{figure}
\end{frame}

\begin{frame}[plain]
    \centering
    GNN Pipeline Explanation
\end{frame}

\begin{frame}{GNN Pipeline}
    \begin{figure}
        \centering
        \includegraphics[width=1\linewidth]{april15th/gnn_architecture.png}
    \end{figure}
\end{frame}

\begin{frame}{MLP Embedding Layer}
    \begin{figure}
        \centering
        \includegraphics[width=0.7\linewidth]{april15th/mlp_embedding_layer_32.png}
    \end{figure}
\end{frame}

\begin{frame}{GNN Pipeline After Embedding}
    \begin{figure}
        \centering
        \includegraphics[width=1\linewidth]{april15th/gnn_architecture_after_embedding.png}
    \end{figure}
\end{frame}

\begin{frame}{GraphConv Layer}
    \begin{figure}
        \centering
        \includegraphics[width=0.8\linewidth]{april15th/graphconv_example.png}
    \end{figure}
\end{frame}

\begin{frame}{GraphConv Layer with Embedding}
    \begin{figure}
        \centering
        \includegraphics[width=0.8\linewidth]{april15th/embedding_gcn_pv.png}
    \end{figure}
\end{frame}

\begin{frame}{Heterogeneous GraphConv Layer with Embedding}
    \begin{figure}
        \centering
        \includegraphics[width=0.8\linewidth]{april15th/heterogeneousgraphconv_embedding.png}
    \end{figure}
\end{frame}

\begin{frame}{GNN Pipeline After GCN}
    \begin{figure}
        \centering
        \includegraphics[width=1\linewidth]{april15th/gnn_architecture_after_gcn(1).png}
    \end{figure}
\end{frame}

\begin{frame}{MLP Edge Predictor Layer}
    \begin{figure}
        \centering
        \includegraphics[width=1\linewidth]{april15th/mlp_edge_predictor_layer.png}
    \end{figure}
\end{frame}

\begin{frame}[plain]
    \centering
    Case Study: Two Overlapping Vertices
\end{frame}

\begin{frame}{Case Study: Two Overlapping Vertices}
    \begin{columns}
        \begin{column}{0.4\textwidth}
            \begin{itemize}
                \item Thr=0.85
                \item Reco PV 26: 59 Assoc tracks
                \begin{itemize}
                    \item \(W_{sum} = 59\)
                    \item Missing Track 168
                    \item Output Probs: 0.834
                \end{itemize}
                \item Truth PV 16: 15 Assoc tracks
                \begin{itemize}
                    \item \(W_{1} = 14\)
                    \item \(W_{1}/W_{sum} = 0.69\)
                \end{itemize}
                \item Truth PV 39: 45 Assoc tracks
                \begin{itemize}
                    \item \(W_{2} = 41\)
                    \item \(W_{2}/W_{sum} = 0.24\)
                \end{itemize}
                \item Classification: Merged
            \end{itemize}
        \end{column}

        \begin{column}{0.6\textwidth}
            \begin{figure}
                \centering
                \includegraphics[width=1\linewidth]{june5th/casestudy_example_reco_truth_hist_dist_16_39.png}
            \end{figure}
        \end{column}
    \end{columns}
\end{frame}

\begin{frame}{Case Study: Two Overlapping Vertices}
    \begin{columns}
        \begin{column}{0.4\textwidth}
            \begin{itemize}
                \item Thr=0.95
                \item Reco PV 26: 51 Assoc tracks
                \begin{itemize}
                    \item \(W_{sum} = 51\)
                \end{itemize}
                \item Truth PV 16: 15 Assoc tracks
                \begin{itemize}
                    \item \(W_{1} = 13\)
                    \item \(W_{1}/W_{sum} = 0.25\)
                \end{itemize}
                \item Truth PV 39: 45 Assoc tracks
                \begin{itemize}
                    \item \(W_{2} = 38\)
                    \item \(W_{2}/W_{sum} = 0.74\)
                \end{itemize}
                \item Classification: Clean
            \end{itemize}
        \end{column}

        \begin{column}{0.6\textwidth}
            \begin{figure}
                \centering
                \includegraphics[width=1\linewidth]{june5th/casestudy_example_reco_truth_hist_dist_16_39.png}
            \end{figure}
        \end{column}
    \end{columns}
\end{frame}

% \begin{frame}{Overview}
%     \centering
%     \textit{Putting it all together!}
%     \begin{itemize}
%         \item PV-Finder Current Studies:
%         \begin{itemize}
%             \item GNN Trained on Truth Vertices
%             \begin{itemize}
%                 \item Nominal/Tight Cuts on Reco Vertex Evaluations
%                 \item Cuts Across Different Thresholds
%             \end{itemize}
%             \item (In Progress) GNN Trained on Reco Vertices
%             \item Hard Scatter Vertex Classification
%             \begin{itemize}
%                 \item Fraction of Events vs Event Type
%                 \item Reco Efficiency vs Local pile-up density
%                 \item TTVA Efficiency and Impurity
%             \end{itemize}
%         \end{itemize}
%         \item Muon Collider Current Studies:
%         \begin{itemize}
%             \item Track Selection Studies
%         \end{itemize}
%     \end{itemize}
% \end{frame}

% \begin{frame}[plain]
%     \centering
%     Training GNN on Truth Vertices
% \end{frame}

% \begin{frame}{Applying Cuts to Track Associations}

% \end{frame}

\begin{frame}[plain]
    \centering
    Training GNN on Reco Vertices
\end{frame}

\begin{frame}{Recap: Steps Outlined}
    \begin{itemize}
        \item Train PV-Finder
        \begin{itemize}
            \item Using current best PV-Finder run (UNet with Dropout Model)
        \end{itemize}
        \item Calculate predicted primary vertex location
        \item Construct graphs from reconstructed PV information
        \item \textbf{Current:} Evaluate using GNN model trained on truth information
        \begin{itemize}
            \item Using current best GNN run (GraphConv Model)
        \end{itemize}
        \item \textbf{New:} Train a GNN model using reco vertices matched to truth vertices
        \begin{itemize}
            \item Allows GNN to learn directly from PV-Finder outputs and associations
        \end{itemize}
    \end{itemize}
\end{frame}

\begin{frame}{Constructing Trainable Data from Reco Vertices}
    \begin{itemize}
        \item Based on Original PV-Finder Classification:
        \begin{itemize}
            \item Clean: Reco tracks (from matched truth vertex) \(\rightarrow\) reco vertex
            \item Merged: Reco tracks (from both matched truth vertex) \(\rightarrow\) reco vertex
            \item Fake: No reco tracks \(\rightarrow\) reco vertex
            \item Split: Remove reco vertices from training
            \begin{itemize}
                \item Don't want GNN to learn features of doubly associating reco tracks
            \end{itemize}
        \end{itemize}
        \item Then you would do vertex classification on test data comparing to truth vertex
        \item \textbf{Work in Progress: Currently training!}
    \end{itemize}
\end{frame}

% \begin{frame}{Reco Vertices with Classification at \(\left< \mu \right> =60\) (Thr=0.95)}
%     \begin{columns}
%         \begin{column}{0.55\textwidth}
%             \begin{itemize}
%                 \item Total Reconstructed PVs:
%                 \item Total Clean: , Rate:
%                 \item Total Merged: , Rate:
%                 \item Total Split: , Rate:
%                 \item Total Fake: , Rate:
%             \end{itemize}
%         \end{column}

%         \begin{column}{0.45\textwidth}
%             \begin{figure}
%                 \centering
%                 \includegraphics[width=1\linewidth]{june5th/reco_pv_avg_bar_chart_thr95_recovertex.png}
%             \end{figure}
%         \end{column}
%     \end{columns}
% \end{frame}

% \begin{frame}{Comparison of Reco Vertex Classification: Thresholds (1)}
%     \begin{columns}
%         \begin{column}{0.5\textwidth}
%             \begin{figure}
%                 \centering
%                 \includegraphics[width=1\linewidth]{june5th/stacked_pv_classification_5_80_thr9_recovertex.png}
%                 \caption{Thr = 0.9}
%             \end{figure}
%         \end{column}

%         \begin{column}{0.5\textwidth}
%             \begin{figure}
%                 \centering
%                 \includegraphics[width=1\linewidth]{june5th/stacked_pv_classification_5_80_thr85_recovertex.png}
%                 \caption{Thr = 0.85}
%             \end{figure}
%         \end{column}
%     \end{columns}
% \end{frame}

% \begin{frame}{Comparison of Reco Vertex Classification: Thresholds (2)}
%     \begin{columns}
%         \begin{column}{0.5\textwidth}
%             \begin{figure}
%                 \centering
%                 \includegraphics[width=1\linewidth]{june5th/stacked_pv_classification_5_80_thr7_recovertex.png}
%                 \caption{Thr = 0.7}
%             \end{figure}
%         \end{column}

%         \begin{column}{0.5\textwidth}
%             \begin{figure}
%                 \centering
%                 \includegraphics[width=1\linewidth]{june5th/stacked_pv_classification_5_80_thr95_recovertex.png}
%                 \caption{Thr = 0.95}
%             \end{figure}
%         \end{column}
%     \end{columns}
% \end{frame}

% \begin{frame}{Number of Associations Per Track vs \( \mu \) Varying Threshold}
%     \begin{columns}
%         \begin{column}{0.5\textwidth}
%             \begin{figure}
%                 \centering
%                 \includegraphics[width=1\linewidth]{june5th/num_track_assoc_vs_mu_thr85.png}
%                 \caption{Thr = 0.85}
%             \end{figure}
%         \end{column}

%         \begin{column}{0.5\textwidth}
%             \begin{figure}
%                 \centering
%                 \includegraphics[width=1\linewidth]{june5th/num_track_assoc_vs_mu_thr95.png}
%                 \caption{Thr = 0.95}
%             \end{figure}
%         \end{column}
%     \end{columns}
% \end{frame}

\begin{frame}[plain]
    \centering
    Hard Scatter Vertex Classification
\end{frame}

\begin{frame}{Methods for Hard Scatter Vertex Classification}
    \begin{itemize}
        \item Usually, primary vertex with highest track \( \sum p_T^2 \) is defined to be HS vertex, \textbf{BUT differs per analysis}
        \begin{itemize}
            \item \( \sum p_T^2 \) methods works for this dataset composed of \(t \bar{t}\) decaying semi-leptonically
        \end{itemize}
        \item Classification Scheme:
        \begin{itemize}
            \item \textbf{Clean/Matched}: \textbf{CLEAN} reco vertex matching true HS interaction AND \textbf{track weights don't contribute \(> 50 \%\) to other vertices}
            \item \textbf{LowPU}: \textbf{MERGED} reco vertex with \textbf{at least 50\%} of the accumulated \textbf{track weight} coming \textbf{from the simulated HS interaction}
            \item \textbf{HighPU}: \textbf{MERGED} reco vertex the HS interaction contributing between \textbf{1\% and 50\% of the accumulated track weight}"
            \item \textbf{PurePU}: No reco vertex with at least 1\% accumulated track weight from HS interaction"
        \end{itemize}
        \item Other ongoing works of interest:
        \begin{itemize}
            \item BDT/GNN-based HS selection: \href{https://indico.cern.ch/event/1374927/contributions/5942568/attachments/2882952/5052279/20240624_hsgnn.pdf}{"Primary vertex selection studies"}
            \begin{itemize}
                \item Related: \href{https://arxiv.org/pdf/2505.19689}{Transformer Jet Flavor Tagging at ATLAS (May 2025)}
            \end{itemize}
            \item \href{https://cds.cern.ch/record/2801638/files/CR2022_018.pdf?version=1}{CMS NN-based PV Reconstruction for L1 Trigger}
        \end{itemize}
    \end{itemize}
\end{frame}

% \begin{frame}{HS Classification: Fraction of Events vs Event Type}

% \end{frame}

% \begin{frame}{HS Classification: Reco Efficiency vs Local pile-up density}

% \end{frame}

% \begin{frame}{Recap: TTVA Efficiency and Impurity}
%     \begin{itemize}
%         \item \textbf{Efficiency: } (\# reco tracks from true HS vertex correctly associated to reco HS vertex) / (\# reco tracks from true HS vertex)
%         \item \textbf{Impurity: } (\# reco tracks from true HS vertex incorrectly associated to reco HS vertex) / (\# reco tracks from reco HS vertex)
%         \item Borrowed from \href{https://indico.cern.ch/event/1025152/contributions/4304219/attachments/2220591/3761247/210406_Basso_ATLAS-Weekly_TTVA-for-Run3.pdf}{Link 1} and \href{https://indico.cern.ch/event/684384/contributions/2805796/attachments/1568829/2473713/uchida_171201.pdf}{Link 2}
%     \end{itemize}
% \end{frame}

% \begin{frame}{HS Vertex: Efficiency and Impurity vs \(\eta\) Thr=0.85}

% \end{frame}

% \begin{frame}{HS Vertex: Efficiency and Impurity vs \(\eta\) Thr=0.95}

% \end{frame}

\begin{frame}{HS Vertex Quality Category Plot}
    \begin{columns}
        \begin{column}{0.5\textwidth}
            \begin{figure}
                \centering
                \includegraphics[width=1\linewidth]{june25th/fig_07_amvf.png}
                \caption{AMVF Pub Note: \href{https://cds.cern.ch/record/2670380/files/ATL-PHYS-PUB-2019-015.pdf}{Link}}
            \end{figure}
        \end{column}

        \begin{column}{0.5\textwidth}
            \begin{figure}
                \centering
                \includegraphics[width=1\linewidth]{june25th/rhs_vertex_quality_category.png}
                \caption{Thr = 0.95}
            \end{figure}
        \end{column}
    \end{columns}
\end{frame}

\begin{frame}{HS Vertex Reconstruction Efficiency}
    \begin{columns}
        \begin{column}{0.3\textwidth}
            \begin{itemize}
                \item Fraction of events which the HS vertex is reconstructed and classified as Clean/Matched, LowPU, or HighPu
                \item Local pile-up density is the number of vertices within a \(\pm\) 2mm window of a HS vertex
            \end{itemize}
        \end{column}

        \begin{column}{0.7\textwidth}
            \begin{figure}
                \centering
                \includegraphics[width=0.85\linewidth]{june25th/vertex_reconstruction_efficiency_vs_local_pu_density.png}
                \caption{Thr = 0.95}
            \end{figure}
        \end{column}
    \end{columns}
\end{frame}

\begin{frame}{Comparison: HS Vertex Reconstruction Efficiency}
    \begin{columns}
        \begin{column}{0.5\textwidth}
            \begin{figure}
                \centering
                \includegraphics[width=1\linewidth]{june25th/fig_08a_amvf_reco.png}
                \caption{AMVF Pub Note (Fig. 8): \href{https://cds.cern.ch/record/2670380/files/ATL-PHYS-PUB-2019-015.pdf}{Link}}
            \end{figure}
        \end{column}

        \begin{column}{0.5\textwidth}
            \begin{figure}
                \centering
                \includegraphics[width=1\linewidth]{june25th/vertex_reconstruction_efficiency_vs_local_pu_density.png}
                \caption{Thr = 0.95}
            \end{figure}
        \end{column}
    \end{columns}
\end{frame}

\begin{frame}{HS Vertex Selection Efficiency}
    \begin{columns}
        \begin{column}{0.3\textwidth}
            \begin{itemize}
                \item Fraction of events where the reconstructed vertex with highest \(\sum p_T^2\) is the one contianing the largest total weight from true HS tracks
            \end{itemize}
        \end{column}

        \begin{column}{0.7\textwidth}
            \begin{figure}
                \centering
                \includegraphics[width=0.85\linewidth]{june25th/vertex_selection_efficiency_vs_local_pu_density.png}
                \caption{Thr = 0.95}
            \end{figure}
        \end{column}
    \end{columns}
\end{frame}

\begin{frame}{Comparison: HS Vertex Selection Efficiency}
    \begin{columns}
        \begin{column}{0.5\textwidth}
            \begin{figure}
                \centering
                \includegraphics[width=1\linewidth]{june25th/fig_09a_amvf_selection.png}
                \caption{AMVF Pub Note (Fig. 9): \href{https://cds.cern.ch/record/2670380/files/ATL-PHYS-PUB-2019-015.pdf}{Link}}
            \end{figure}
        \end{column}

        \begin{column}{0.5\textwidth}
            \begin{figure}
                \centering
                \includegraphics[width=1\linewidth]{june25th/vertex_selection_efficiency_vs_local_pu_density.png}
                \caption{Thr = 0.95}
            \end{figure}
        \end{column}
    \end{columns}
\end{frame}

\backupend

\end{document}
